% preamble-exam-zh.tex — 共享 preamble
% 用于所有 exam-zh 模板
%
% 样式策略：
%   屏幕 → 语义彩色（蓝=关键想法，红=易错，绿=路线，灰=提示/教师备注）
%   打印 → 自动降级为黑白（通过 print mode 或 monochrome 选项）

\documentclass{exam-zh}

% ---------- 基础配置 ----------
\examsetup{
  page/size = a4paper,
  page/show-head = false,
  page/show-foot = false,
  sealline/show = false,
  font = times,
  math-font = xits,
}

\usepackage{tcolorbox}
\usepackage{needspace}
\tcbuselibrary{skins,breakable}

% ---------- 打印/屏幕颜色切换 ----------
% 用法：\documentclass[monochrome]{exam-zh} 即可切换为纯黑白
% 注意：不要使用 \if@xxx 放在 makeatother 外部；这里改成普通条件名。
\newif\ifeduprintcolor
\eduprintcolortrue

\makeatletter
\@ifclasswith{exam-zh}{monochrome}{\eduprintcolorfalse}{}
\makeatother

% ---------- 语义颜色定义 ----------
% 屏幕：有色彩区分；打印：降级为灰度
\ifeduprintcolor
  % --- 屏幕彩色模式 ---
  \definecolor{edu-blue}{RGB}{47, 82, 122}       % 关键想法
  \definecolor{edu-blue-bg}{RGB}{243, 246, 252}
  \definecolor{edu-red}{RGB}{180, 40, 40}         % 易错提醒
  \definecolor{edu-red-bg}{RGB}{255, 245, 240}
  \definecolor{edu-green}{RGB}{34, 100, 60}       % 路线图
  \definecolor{edu-green-bg}{RGB}{242, 250, 245}
  \definecolor{edu-orange}{RGB}{160, 90, 0}       % 训练目标
  \definecolor{edu-orange-bg}{RGB}{255, 248, 235}
  \definecolor{edu-gray}{RGB}{100, 100, 100}      % 提示/教师备注
  \definecolor{edu-gray-bg}{RGB}{248, 248, 248}
  \definecolor{edu-purple}{RGB}{90, 50, 120}      % 思考题
  \definecolor{edu-purple-bg}{RGB}{248, 244, 252}
\else
  % --- 打印黑白模式 ---
  \definecolor{edu-blue}{gray}{0.15}
  \definecolor{edu-blue-bg}{gray}{0.96}
  \definecolor{edu-red}{gray}{0.25}
  \definecolor{edu-red-bg}{gray}{0.97}
  \definecolor{edu-green}{gray}{0.20}
  \definecolor{edu-green-bg}{gray}{0.97}
  \definecolor{edu-orange}{gray}{0.25}
  \definecolor{edu-orange-bg}{gray}{0.98}
  \definecolor{edu-gray}{gray}{0.40}
  \definecolor{edu-gray-bg}{gray}{0.97}
  \definecolor{edu-purple}{gray}{0.30}
  \definecolor{edu-purple-bg}{gray}{0.98}
\fi

% ---------- 自定义教学环境 ----------
% 共性：enhanced + breakable（跨页不断裂）

% 原题卡片 — 强边框，突出显示
\newtcolorbox{problemcard}{
  enhanced, breakable,
  colback=edu-blue-bg,
  colframe=edu-blue,
  boxrule=1.2pt,
  arc=0pt,
  left=3mm, right=3mm, top=3mm, bottom=3mm,
}

% 解题路线图 — 左边线 + 浅底
\newtcolorbox{routemap}{
  enhanced, breakable,
  colback=edu-green-bg,
  colframe=edu-green!30,
  boxrule=0pt,
  borderline west={2.5pt}{0pt}{edu-green},
  left=3mm, right=3mm, top=3mm, bottom=3mm,
}

% 关键想法 — 蓝色左边线
\newtcolorbox{keyideabox}{
  enhanced, breakable,
  colback=edu-blue-bg,
  colframe=edu-blue!30,
  boxrule=0pt,
  borderline west={2.5pt}{0pt}{edu-blue},
  left=3mm, right=3mm, top=3mm, bottom=3mm,
}

% 易错提醒 — 红色左边线
\newtcolorbox{mistakebox}{
  enhanced, breakable,
  colback=edu-red-bg,
  colframe=edu-red!20,
  boxrule=0pt,
  borderline west={3pt}{0pt}{edu-red},
  left=3mm, right=3mm, top=2.5mm, bottom=2.5mm,
}

% 提示 — 灰色虚线左边线
\newtcolorbox{hintbox}{
  enhanced, breakable,
  colback=edu-gray-bg,
  colframe=edu-gray!20,
  boxrule=0pt,
  borderline west={2pt}{0pt}{edu-gray, dashed},
  left=3mm, right=3mm, top=2mm, bottom=2mm,
}

% 思考题 — 紫色虚线左边线
\newtcolorbox{thinkbox}{
  enhanced, breakable,
  colback=edu-purple-bg,
  colframe=edu-purple!20,
  boxrule=0pt,
  borderline west={2pt}{0pt}{edu-purple, dashed},
  left=2.5mm, right=2.5mm, top=2.5mm, bottom=2.5mm,
}

% 教师备注 — 灰色左边线，小字
\newtcolorbox{teachernote}{
  enhanced, breakable,
  colback=edu-gray-bg,
  colframe=edu-gray!15,
  boxrule=0pt,
  borderline west={2pt}{0pt}{edu-gray!60},
  left=2.5mm, right=2.5mm, top=2.5mm, bottom=2.5mm,
  fontupper=\small,
  colupper=black!60,
}

% 训练目标 — 橙色左边线，小字
\newtcolorbox{traininggoal}{
  enhanced, breakable,
  colback=edu-orange-bg,
  colframe=edu-orange!15,
  boxrule=0pt,
  borderline west={1.5pt}{0pt}{edu-orange},
  left=2mm, right=2mm, top=1mm, bottom=1mm,
  fontupper=\small,
}

% 答题区 — 无边框，纯留白
\newcommand{\answerarea}[1][40mm]{%
  \vspace{2mm}%
  \begin{tcolorbox}[
    enhanced, breakable,
    colback=white,
    colframe=white,
    boxrule=0pt,
    height=#1,
    valign=top,
    bottom=0mm,
    after skip=0mm,
  ]
  \end{tcolorbox}%
}

% ---------- 字体设置 ----------
% exam-zh (ctexbook) already sets CJK fonts via ctex.
% Only override if specific fonts are needed.
% macOS: Songti SC, STHeiti, STFangsong
% Linux: SimSun, SimHei, FangSong

% ---------- 页面设置 ----------
% exam-zh (ctexbook) already loads geometry.
% Override margins if needed:
% \geometry{top=16mm, bottom=16mm, left=14mm, right=14mm}

\examsetup{
  question/show-answer = true,
  fillin/show-answer = true,
  solution/show-solution = show-stay,
}

% ---------- 教师版解答样式 ----------
\newtcolorbox{solutionblock}{
  enhanced, breakable,
  colback=edu-blue-bg,
  colframe=edu-blue!25,
  boxrule=0pt,
  borderline west={2pt}{0pt}{edu-blue!50},
  arc=0pt,
  left=3mm, right=3mm, top=2mm, bottom=2mm,
  before skip=2mm,
  after skip=3mm,
  fontupper=\small,
}

% 解答步骤编号
\newcounter{solstep}
\newcommand{\solstep}[2]{%
  \stepcounter{solstep}%
  \par\medskip
  \noindent\hangindent=6mm
  {\color{edu-blue}\bfseries\textcircled{\scriptsize\thesolstep}\enspace #1}\enspace #2\par
}

\begin{document}

\section*{求点坐标基础}

\section*{一、选择题（每题 12 分，共 48 分）}



\begin{keyideabox}
  \textbf{核心思想：坐标与方程的统一性}
  几何上的‘点在直线上’等价于代数上的‘点坐标满足解析式’。求交点本质上是寻找同时满足两个方程的公共解。
\par\small 数形结合是解析几何的根本方法。\end{keyideabox}



\begin{question}[points=12]
  直线 $y = kx - 5$ 经过点 $(3, 1)$，则 $k$ 的值为
  \begin{choices}
    \item $2$
    \item $-2$
    \item $\dfrac{5}{3}$
    \item $-\dfrac{5}{3}$
  \end{choices}
  \paren[A]
  \begin{solutionblock}
    代入 $(3, 1)$：$1=3k-5$，$3k=6$，$k=2$。
  \end{solutionblock}
\end{question}



\begin{question}[points=12]
  直线 $y = -x + 5$ 与坐标轴的交点分别为
  \begin{choices}
    \item $(5, 0)$ 和 $(0, 5)$
    \item $(-5, 0)$ 和 $(0, 5)$
    \item $(5, 0)$ 和 $(0, -5)$
    \item $(-5, 0)$ 和 $(0, -5)$
  \end{choices}
  \paren[A]
  \begin{solutionblock}
    令 $y=0$：$x=5$，$A(5,0)$。令 $x=0$：$y=5$，$B(0,5)$。
  \end{solutionblock}
\end{question}



\begin{question}[points=12]
  直线 $y = 3x + 2$ 与 $y = -x + 6$ 的交点坐标为
  \begin{choices}
    \item $(1, 1)$
    \item $(2, 5)$
    \item $(1, 5)$
    \item $(2, 4)$
  \end{choices}
  \paren[C]
  \begin{solutionblock}
    令 $3x+2=-x+6$，$4x=4$，$x=1$。$y=3+2=5$。交点 $(1, 5)$。
  \end{solutionblock}
\end{question}



\begin{question}[points=12]
  直线 $y = \dfrac{1}{2}x - 4$ 与 $x$ 轴的交点为
  \begin{choices}
    \item $(8, 0)$
    \item $(-8, 0)$
    \item $(0, -4)$
    \item $(2, 0)$
  \end{choices}
  \paren[A]
  \begin{solutionblock}
    令 $y=0$：$0=\frac{1}{2}x-4$，$x=8$。
  \end{solutionblock}
\end{question}


\section*{二、填空题（每题 12 分，共 12 分）}



\begin{question}[points=12]
  直线 $y = 3x + 3$ 与 $y = -2x + 8$ 的交点纵坐标为 \rule{3cm}{0.4pt}。
  \fillin
  \paren[$6$]
  \begin{solutionblock}
    令 $3x+3=-2x+8$，$5x=5$，$x=1$，$y=3+3=6$。
  \end{solutionblock}
\end{question}


\section*{三、解答题（共 40 分）}




\needspace{8\baselineskip}

\begin{problem}[points=20]
  直线 $y = kx + 8$ 经过点 $A(-2, 0)$。

(1) 求 $k$ 的值；
(2) 写出直线解析式，并求该直线与 $y$ 轴的交点坐标。
  \answerarea[60mm]
  \setcounter{solstep}{0}
  \begin{solutionblock}
    \textbf{答案：}(1) $k=4$；(2) $y=4x+8$，$y$ 轴交点 $(0, 8)$\par\medskip
    \solstep{第(1)问}{代入 $(-2, 0)$：$0 = -2k + 8$，$2k = 8$，$k = 4$。}
    \solstep{第(2)问}{解析式 $y = 4x + 8$。令 $x=0$：$y=8$。$y$ 轴交点为 $(0, 8)$。}
  \end{solutionblock}
\end{problem}


\needspace{8\baselineskip}

\begin{problem}[points=20]
  直线 $l_1$：$y = -x + 5$ 与直线 $l_2$：$y = 2x - 1$。

求 $l_1$ 与 $l_2$ 的交点坐标。
  \answerarea[50mm]
  \setcounter{solstep}{0}
  \begin{solutionblock}
    \textbf{答案：}$(2, 3)$\par\medskip
    \solstep{联立}{$-x + 5 = 2x - 1$}
    \solstep{解 $x$}{$6 = 3x$，$x = 2$}
    \solstep{回代}{代入 $l_1$：$y = -2 + 5 = 3$。交点 $(2, 3)$。}
  \end{solutionblock}
\end{problem}



\end{document}
